\section{Tanda Tangan Digital}

Tanda tangan digital memastikan keaslian dan integritas dokumen digital \cite{rivest1978,katz2020,ferguson2010}. Prinsip: pengirim mengenkripsi hash dokumen dengan kunci privat; penerima mendekripsi dengan kunci publik pengirim dan memverifikasi hash cocok dengan hash dokumen yang dihitung ulang.

\begin{figure}[!htbp]
\centering
\begin{tikzpicture}[
  node distance=1cm,
  box/.style={rectangle, draw, minimum width=2cm, minimum height=0.6cm, align=center, font=\small}
]
\node[box, fill=blue!15] (alice) {Alice};
\node[box, fill=green!15, right=3.5cm of alice] (bob) {Bob};
\draw[-Stealth] (alice) -- node[above, font=\scriptsize] {$m, s = H(m)^d$} (bob);
\node[below=0.5cm of alice, font=\scriptsize] {Kunci privat};
\node[below=0.5cm of bob, font=\scriptsize] {Verifikasi: $s^e \stackrel{?}{=} H(m)$};
\end{tikzpicture}
\caption{Alur tanda tangan digital berbasis RSA}
\label{fig:tanda-tangan}
\end{figure}

Skema berbasis RSA: Sign$(m) = H(m)^d \bmod n$; Verify$(m, s) = (s^e \bmod n \stackrel{?}{=} H(m))$. Tanda tangan digital memberikan non-repudiation: pengirim tidak dapat menyangkal telah menandatangani \cite{rivest1978,katz2020,ferguson2010}.
